Nowadays, Internet of Things (IoT) devices and technologies are getting more and more importance in our everyday lives and are almost everywhere. 
The motivations behind their fast development are easy to understand: 
first of all, Micro-Electro-Mechanical Systems (MEMS) sensors are small, easy to be integrated and are getting cheaper; 
modern microcontrollers are easy to get and program and, 
finally, recent evolution of batteries allowed the development of IoT nodes with appropriate lifetime. 
Hence, these small devices capable of sending sensor data for further processing are getting popularity in different fields. 
Despite all of this, these systems present their drawbacks: data coming from these MEMS sensors is heavily affected by noise and biases, low-energy transmission protocols present low throughput and significant delays, the computational power of the microcontrollers is limited and a large number of factors affect the battery duration of these devices.
Developers are forced to find compromises between measurement accuracy, battery duration and system costs, trying to optimize as much as possible all the variables coming into play. 

An interesting application for these sensor systems is sensor fusion, which consists in joining together heterogeneous data coming from different sensors to get accurate information. 
Different algorithms are available, and often some are already available in the firmware of sensor boards from different manufacturers. 
For example, ST Microelectronics provides Human Activity Recognition algorithms in their evaluation boards, and posture recognition for health purposes is another possible application. 
To generate satisfying results, these algorithms need the sensor signals to be synchronized in time: this is a big challenge when dealing with IoT and embedded systems, and it is even more important in critical applications. 
For example, when considering object-recognition algorithms in car accident-prevention systems, accuracy is heavily impacted when data from camera (2D information) and LiDAR (depth) is not correctly synchronized. 
The literature presents some works on clock alignment, but those are using low-layer APIs, are intended for the same HW platform and aren't always thought for wireless transmission. Hence, even if the presented results are promising, those methods can’t be easily applied to heterogeneous, mass-produced devices. 

In this thesis I present the system developed to synchronize the wireless-based acquisition of signals from four different sensor boards from ST Microelectronics and three from MBIENTLAB.
I also explored the possibility to use the embedded microphones and the consequent audio data to synchronize and label data, finding valid ways to stream sound and implement Fourier Fast Transform (FFT) on the microcontrollers. 
I gave particular attention in getting the most out of the Bluetooth Low Energy protocol (BLE): it is a good solution from the energy-efficiency and compatibility points of view, but at the same time some optimization is necessary to keep a good throughput and maintain an efficient bi-directional connection. 
Finally, I developed time synchronization on the application layer of the BLE stack: it consists in the periodic exchange of timestamped packets between the client and the sensors. 
The client is then aware of the internal timestamp of the sensors in those re-synchronization moment, and can tag the samples with the precise moment in time they are acquired. 
The goals are 
a) development of firmware for the open-source ST devices, with a personalized time synchronization mechanism, 
b) design of clients for all the devices, formatting data in a consistent way and 
c) writing of MATLAB scripts to join together the resulting data. 

The developed system behaves as expected in the experiments that have been carried out. 
I made initial tests to determine the best data rate for the sensor information (to keep a packet loss/received ratio lower than 1\%), which is 75Hz. 
Then, I verified the behavior of the time synchronization after 15 minutes of data streaming, showing a precision of around 70 ms, while the same setup without synchronization presents samples with delays of more than 1s. 
Finally, I performed an extensive test by installing the sensors in a car and collecting data in that environment, with the help of a camera and GPS devices. 
In this case, data have been labelled using sounds at predefined frequencies and elaborated at the node level. 
Results demonstrated that sensor information remains aligned in time and consistent with the recorded audio data.
